\documentclass[11pt,a4paper]{article}
\usepackage[utf8]{inputenc}
\usepackage[T1]{fontenc}
\usepackage{times}
\usepackage[margin=2.5cm]{geometry}
\usepackage{booktabs}
\usepackage{amsmath,amssymb}
\usepackage{graphicx}
\usepackage{xcolor}
\usepackage{hyperref}
\usepackage{multirow}
\usepackage{enumitem}
\usepackage{tabularx}
\usepackage[font=small]{caption}
\usepackage{subcaption}

\hypersetup{colorlinks=true,linkcolor=blue!60!black,citecolor=blue!60!black,urlcolor=blue!60!black}

\title{Differential Depth Response: Fingerprinting Generation Depth\\in Non-Paraphrastic LLM Chains}

\author{
  Anonymous Submission to EMNLP 2026
}
\date{}

\begin{document}
\maketitle

%% ============================================================
%% ABSTRACT
%% ============================================================
\begin{abstract}
Web-scraped training corpora increasingly contain text produced by large language models (LLMs), yet current detection methods offer only binary human-vs-synthetic classification, conflating qualitatively different material across generation depths.
We formulate \emph{generation depth classification}---distinguishing human-authored text (gen-0), first-generation synthetic text (gen-1), and recursively generated text (gen-2+)---and characterize a phenomenon we call \emph{differential depth response} (DDR): surface-level diversity features (hapax ratio, type--token ratio) exhibit strong monotonic drift with depth ($|\rho| > 0.55$), entropy-based features show moderate drift ($|\rho| \in [0.18, 0.52]$), while distributional shape features (tail decay, entropy max) remain near-stationary ($|\rho| < 0.10$).
A lightweight 7-feature ensemble achieves leak-free binary F1 of 0.675 under leave-one-generator-out (LOGO) evaluation across 5~generators from 3~model families (0.6B--2.8B), outperforming the best single-detector baseline by 2.2 percentage points ($p = 0.008$).
The signal persists across domains (F1 range 0.648--0.692) and prompt types, but degrades at larger model scales (8B: entropy mean $|\rho|$ drops to 0.36).
A downstream fine-tuning experiment confirms that depth-aware data curation does not improve language modeling quality (TOST $p = 0.00012$ at $\pm 2\%$ margin), positioning depth fingerprinting as a \emph{diagnostic} tool for provenance characterization rather than a curative data-selection mechanism.
\end{abstract}

%% ============================================================
%% 1. INTRODUCTION
%% ============================================================
\section{Introduction}

Recent estimates suggest that 7--17\% of scientific publications and an increasing fraction of web text are partially or fully generated by large language models \cite{liang2024mapping}.
As LLM-generated text enters training corpora, it may be recursively consumed by subsequent models---a feedback loop termed \emph{model collapse} in the theoretical literature \cite{shumailov2024model}.
Characterizing the \emph{generation depth} of text---how many recursive generation steps separate it from human authorship---is therefore a pressing concern for data governance, provenance auditing, and understanding distributional shifts in training data.

Existing detection methods frame the problem as binary classification: human vs.\ machine \cite{mitchell2023detectgpt,hans2024binoculars,bao2024fastdetectgpt}.
This binary framing conflates qualitatively different material: first-generation text (gen-1) retains much of the source distribution's lexical properties, while deeper text exhibits measurably reduced vocabulary diversity and compressed perplexity \cite{briesch2023large,guo2024curious}.
Moreover, paraphrastic chains converge to attractor cycles within 2--3 iterations \cite{chen2025llm}, bounding the discriminability of semantic-preserving rewrites.
No existing method addresses generation depth estimation in non-paraphrastic continuation chains, where the absence of semantic-preservation constraints allows distributional drift to accumulate over deeper generation steps.

We formulate \textbf{generation depth classification} as a three-class ordinal task (gen-0 / gen-1 / gen-2+) and introduce the concept of \textbf{differential depth response} (DDR): the empirical observation that distributional features respond to generation depth at different rates.
Rather than a binary drifting/stationary dichotomy, we find a \emph{three-gradient} structure: \textbf{strong drift} features (hapax ratio $\rho = {-}0.61$, TTR $\rho = {-}0.58$), \textbf{moderate drift} features (entropy mean $\rho = {-}0.52$, entropy min $\rho = {-}0.31$, entropy std $\rho = {-}0.18$), and \textbf{near-stationary} features (entropy max $\rho = 0.08$, tail decay $\rho = 0.05$).
This gradient structure is consistent with recursive generation inducing a location shift in the next-token log-likelihood (NLL) distribution while approximately preserving distributional shape.

Our contributions are:
\begin{enumerate}[nosep,leftmargin=*]
  \item \textbf{Task formulation:} We define generation depth classification beyond binary detection, with a leak-free evaluation protocol (LOGO) that prevents seed-sharing inflation.
  \item \textbf{Empirical characterization of DDR:} We document the three-gradient DDR structure across 5 generators from 3 model families (0.6B--2.8B), showing it is robust across reference models, domains, and prompt types, but degrades at 8B scale.
  \item \textbf{Negative finding:} Depth-aware data curation is statistically equivalent to random subsampling for downstream language modeling quality, establishing depth fingerprinting as a diagnostic rather than curative tool.
\end{enumerate}

%% ============================================================
%% 2. RELATED WORK
%% ============================================================
\section{Related Work}

\paragraph{LLM-Generated Text Detection.}
Statistical detection methods leverage distributional signatures of machine-generated text.
DetectGPT \cite{mitchell2023detectgpt} uses perturbation-based log-probability curvature; Fast-DetectGPT \cite{bao2024fastdetectgpt} replaces perturbation with conditional sampling for efficiency.
Binoculars \cite{hans2024binoculars} computes cross-perplexity ratios between two models.
GLTR \cite{gehrmann2019gltr} visualizes per-token rank distributions.
DNA-GPT \cite{yang2024dna} exploits divergent n-gram patterns under re-generation.
All operate as binary classifiers and do not estimate generation depth.

\paragraph{Model Collapse and Iterative Generation.}
Shumailov et al.\ \cite{shumailov2024model} demonstrate that recursive training on synthetic data causes distributional tails to vanish---\emph{model collapse}.
Briesch et al.\ \cite{briesch2023large} show vocabulary diversity decreases across generation iterations.
Guo et al.\ \cite{guo2024curious} characterize perplexity compression in iterative text generation.
Our work differs in treating iterative generation as a \emph{classification signal} rather than a failure mode to prevent.

\paragraph{Paraphrastic Chains and Attractor Theory.}
Chen et al.\ \cite{chen2025llm} model LLM paraphrasing as a Markov chain converging to attractor cycles, bounding discriminability to $\sim$3 classes.
Their semantic-preservation constraint limits diversity drift; our non-paraphrastic continuation chains lack this constraint, enabling depth signals to persist beyond depth 4.

\paragraph{Data Quality and Provenance.}
Penedo et al.\ \cite{penedo2024fineweb} score web documents by educational value.
Watermarking methods \cite{kirchenbauer2023watermark} embed provenance markers but require generator cooperation.
Depth fingerprinting is complementary: it operates post-hoc without generator access.

%% ============================================================
%% 3. METHOD
%% ============================================================
\section{Method}

\subsection{Task Formulation}

Given a text passage $x$, we define the \emph{generation depth} $d(x) \in \{0, 1, 2, 3, \ldots\}$ as the number of recursive LLM generation steps separating $x$ from its human-authored seed.
We formulate three classification granularities:
\textbf{binary} (gen-1 vs.\ gen-2+, excluding gen-0 to avoid seed-sharing leakage),
\textbf{3-class} (gen-0 / gen-1 / gen-2+), and
\textbf{5-class} (gen-0 through gen-4).
Our primary metric is \textbf{leak-free binary F1} under leave-one-generator-out (LOGO) evaluation, which prevents both seed-sharing and generator memorization.

\subsection{Generation Chains}

We construct non-paraphrastic continuation chains.
Starting from human-authored seed texts drawn from Wikipedia, news, scientific, and creative writing corpora (500 seeds per domain), we generate depth-1 text by prompting a generator model to continue the seed.
Each subsequent depth uses the previous generation as input, producing chains of depth 0--4.
We use 5 generators: Pythia-2.8B, Pythia-1B, OLMo-2-1B, Qwen3-1.7B, and Qwen3-0.6B, yielding 22,800 total text samples.
All generation uses temperature $t = 0.7$.

\subsection{Feature Extraction}

We extract 7 features from each text, computed relative to a reference model (Pythia-410M by default):

\paragraph{Surface features (2).}
\emph{Type--token ratio} (TTR) and \emph{hapax ratio} (proportion of words appearing exactly once) capture vocabulary diversity directly from the text without model access.

\paragraph{Reference-based features (5).}
We compute the per-token NLL distribution under the reference model and extract five statistics: \emph{entropy mean}, \emph{entropy std}, \emph{entropy min}, \emph{entropy max}, and \emph{tail decay} (rate parameter of exponential fit to the upper tail of the NLL distribution).

\subsection{Classifier}

We use XGBoost \cite{chen2016xgboost} with default hyperparameters as our classifier, chosen for interpretability and efficiency.
For cross-generator evaluation, we employ \textbf{LOGO}: train on 4 generators, test on the held-out generator, repeating across all 5 folds.
We also evaluate leave-one-family-out (\textbf{LOFO}), where entire model families are held out.

%% ============================================================
%% 4. EXPERIMENTS
%% ============================================================
\section{Experiments}

\subsection{DDR Analysis and Feature Ablation}
\label{sec:ddr}

We first characterize the differential depth response across our 7 features.
Table~\ref{tab:ddr_gradient} shows Spearman rank correlations ($\rho$) between each feature and generation depth.
Features naturally partition into three gradient groups rather than the binary drifting/stationary split one might expect.

\begin{table}[t]
\centering
\small
\caption{DDR gradient classification. Features exhibit three distinct response levels to generation depth.}
\label{tab:ddr_gradient}
\begin{tabular}{lcc}
\toprule
\textbf{Feature} & \textbf{Spearman $\rho$} & \textbf{Gradient} \\
\midrule
hapax ratio & $-$0.61 & Strong drift \\
TTR & $-$0.58 & Strong drift \\
\midrule
entropy mean & $-$0.52 & Moderate drift \\
entropy min & $-$0.31 & Moderate drift \\
entropy std & $-$0.18 & Moderate drift \\
\midrule
entropy max & $+$0.08 & Near-stationary \\
tail decay & $+$0.05 & Near-stationary \\
\bottomrule
\end{tabular}
\end{table}

Surface features (hapax ratio, TTR) show the strongest monotonic drift ($|\rho| > 0.55$), reflecting progressive vocabulary impoverishment as text is recursively generated.
Entropy-based features occupy a moderate range: entropy mean ($\rho = {-}0.52$) captures the location shift in the NLL distribution, while entropy std ($\rho = {-}0.18$) and entropy min ($\rho = {-}0.31$) reflect secondary distributional effects.
Shape features (entropy max, tail decay) remain near-stationary, consistent with the NLL distribution undergoing a location shift that preserves upper-tail structure.

Table~\ref{tab:ablation} presents feature ablation results under 5-fold cross-validation.
The full 7-feature set (macro F1 = 0.808) outperforms both surface-only (0.757) and reference-only (0.712) subsets, confirming feature complementarity (McNemar $p \approx 0$ for both comparisons).
Surface-only outperforms reference-only by 4.5~pp, indicating that vocabulary diversity is the dominant signal source.

\begin{table}[t]
\centering
\small
\caption{Feature ablation (5-fold CV, macro F1). Full feature set outperforms both subsets.}
\label{tab:ablation}
\begin{tabular}{lcc}
\toprule
\textbf{Feature Set} & \textbf{Macro F1} & \textbf{95\% CI} \\
\midrule
Full (7 features) & 0.808 & {[0.795, 0.821]} \\
Surface-only (TTR + hapax) & 0.757 & {[0.742, 0.772]} \\
Reference-only (5 entropy + tail) & 0.712 & {[0.698, 0.726]} \\
\bottomrule
\end{tabular}
\end{table}

Median perplexity decreases monotonically with depth (13.2, 5.1, 3.1, 2.2, 1.8 for depths 0--4), validating the premise that deeper generations become progressively more predictable under the reference model.

\subsection{Cross-Generator Generalization}
\label{sec:cross_gen}

Table~\ref{tab:logo} reports LOGO evaluation across 5 generators.
The average LOGO binary F1 is 0.675 (95\% CI [0.662, 0.688]), with per-generator scores ranging from 0.620 (Qwen3-0.6B) to 0.715 (Pythia-1B)---a 9.5~pp spread.
LOFO binary F1 is 0.652, reflecting a 2.3~pp drop when entire model families are held out.

\begin{table}[t]
\centering
\small
\caption{LOGO per-generator binary F1. Pythia generators benefit from reference-model family proximity.}
\label{tab:logo}
\begin{tabular}{lccc}
\toprule
\textbf{Generator} & \textbf{Binary F1} & \textbf{95\% CI} & \textbf{Family} \\
\midrule
Pythia-1B & 0.715 & {[0.697, 0.733]} & Pythia \\
Pythia-2.8B & 0.710 & {[0.692, 0.728]} & Pythia \\
Qwen3-1.7B & 0.668 & {[0.649, 0.687]} & Qwen \\
OLMo-2-1B & 0.660 & {[0.641, 0.679]} & OLMo \\
Qwen3-0.6B & 0.620 & {[0.600, 0.640]} & Qwen \\
\midrule
\textbf{LOGO avg} & \textbf{0.675} & {[0.662, 0.688]} & --- \\
\textbf{LOFO} & \textbf{0.652} & --- & --- \\
\bottomrule
\end{tabular}
\end{table}

The Pythia generators (2.8B, 1B) achieve higher F1 due to family proximity with the Pythia-410M reference model.
Qwen3-0.6B scores lowest, likely because its limited model capacity constrains generation diversity, compressing depth drift signals.
The LOFO degradation (2.3~pp) indicates mild cross-family transfer loss, primarily from removing the two Pythia generators that benefit from reference-model affinity.

\subsection{Comparison with Detection Baselines}
\label{sec:baselines}

We evaluate existing detection scores as ordinal depth predictors by training separate XGBoost classifiers on each single-feature score under the same LOGO protocol.
Table~\ref{tab:baselines} shows that the 7-feature ensemble (F1 = 0.680) outperforms the best single-detector baseline, avg-logprob (F1 = 0.658), by 2.2~pp (bootstrap $p = 0.008$).

\begin{table}[t]
\centering
\small
\caption{LOGO binary F1: 7-feature ensemble vs.\ single-detector baselines. Advantage is statistically significant but modest.}
\label{tab:baselines}
\begin{tabular}{lccc}
\toprule
\textbf{Method} & \textbf{LOGO F1} & \textbf{95\% CI} & \textbf{$\rho$} \\
\midrule
7-feature ensemble & 0.680 & {[0.665, 0.695]} & --- \\
Avg-LogProb & 0.658 & {[0.540, 0.740]} & $-$0.52 \\
LD-Score & 0.632 & {[0.520, 0.710]} & $-$0.48 \\
FastDetectGPT & 0.595 & {[0.480, 0.690]} & $-$0.38 \\
Binoculars & 0.558 & {[0.400, 0.650]} & $-$0.20 \\
\bottomrule
\end{tabular}
\end{table}

The 7-feature advantage derives from combining vocabulary diversity (surface features) with distributional statistics, providing a more robust decision boundary than any single detector score.
However, we emphasize that the 2.2~pp margin is modest; the contribution lies not in the performance gap itself but in (a) the task formulation, (b) the DDR characterization, and (c) the interpretability and efficiency of the feature-based approach.
Notably, single-detector baselines exhibit wide confidence intervals (reflecting high cross-generator variance), making them unreliable ordinal predictors in isolation.

\subsection{Cross-Model Chains}
\label{sec:cross_model}

Real-world recursive generation may involve different models at each depth step.
We construct cross-model chains by randomly assigning one of 5 generators to each depth step.
Table~\ref{tab:cross_model} compares cross-model and single-model chain results.

\begin{table}[t]
\centering
\small
\caption{Cross-model vs.\ single-model chains. Architecture mixing provides moderate amplification of depth signal.}
\label{tab:cross_model}
\begin{tabular}{lcccc}
\toprule
\textbf{Chain Type} & \textbf{Bin.\ F1} & \textbf{3-cl.\ F1} & \textbf{hapax $\rho$} & \textbf{ent.\ mean $\rho$} \\
\midrule
Cross-model & 0.730 & 0.720 & $-$0.54 & $-$0.40 \\
Single-model (LOGO) & 0.675 & 0.788 & $-$0.61 & $-$0.52 \\
\bottomrule
\end{tabular}
\end{table}

Cross-model chains yield higher binary F1 (0.730 vs.\ 0.675, $\Delta = 5.5$~pp), as architecture differences between generators introduce additional distributional shifts that compound with depth-intrinsic drift.
However, this amplification is moderate---not sufficient to support a strong ``heterogeneity amplification'' narrative.
The reduced $|\rho|$ for entropy mean ($-0.40$ vs.\ $-0.52$) reflects inconsistent entropy estimation across different architectures, while surface features (hapax $\rho$) degrade less ($-0.54$ vs.\ $-0.61$), confirming their relative robustness to architecture variation.

\subsection{Scale Effect}
\label{sec:scale}

To test whether DDR persists at larger scales, we add Qwen3-8B-Base to our generator set (Table~\ref{tab:scale}).

\begin{table}[t]
\centering
\small
\caption{Scale effect on DDR. Larger models compress entropy-based drift while surface features remain relatively stable.}
\label{tab:scale}
\begin{tabular}{lcccc}
\toprule
\textbf{Generator} & \textbf{Size} & \textbf{Bin.\ F1} & \textbf{hapax $\rho$} & \textbf{ent.\ mean $\rho$} \\
\midrule
Qwen3-0.6B & 0.6B & 0.620 & $-$0.58 & $-$0.52 \\
Pythia-1B & 1B & 0.715 & $-$0.63 & $-$0.54 \\
OLMo-2-1B & 1.5B & 0.660 & $-$0.60 & $-$0.50 \\
Qwen3-1.7B & 1.7B & 0.668 & $-$0.59 & $-$0.49 \\
Pythia-2.8B & 2.8B & 0.710 & $-$0.61 & $-$0.52 \\
\midrule
Qwen3-8B & 8B & 0.615 & $-$0.50 & $-$0.36 \\
\bottomrule
\end{tabular}
\end{table}

Qwen3-8B shows a marked reduction in entropy mean drift ($|\rho| = 0.36$ vs.\ the 0.49--0.54 range of smaller models), while surface feature drift degrades more gradually (hapax $|\rho| = 0.50$ vs.\ 0.58--0.63).
This asymmetric degradation arises because larger models produce more fluent text with narrower next-token entropy distributions, compressing the NLL location shift that drives entropy-based DDR.
Vocabulary diversity compression, however, is an intrinsic property of recursive generation that persists regardless of model capacity---hapax and TTR continue to drift because each generation step reuses and reinforces the vocabulary of its input.

Including Qwen3-8B in the LOGO pool slightly lowers the overall average (0.672 vs.\ 0.675 for 5 generators).
This scale-dependent degradation is a key limitation: at frontier model scales ($> 14$B), entropy features may shift from moderate to near-stationary drift, and surface features would need to carry the classification signal alone.

\subsection{Downstream Utility}
\label{sec:downstream}

We test whether depth fingerprinting can improve downstream language modeling through depth-aware data curation.
We fine-tune OLMo-1B on depth-stratified training data (3 seeds, bootstrap CIs).
Table~\ref{tab:downstream} reports perplexity on a held-out evaluation set.

\begin{table}[t]
\centering
\small
\caption{Downstream CLM quality by training data composition. Depth-aware curation does not outperform random subsampling.}
\label{tab:downstream}
\begin{tabular}{lccc}
\toprule
\textbf{Variant} & \textbf{Train Size} & \textbf{Ppl} & \textbf{TOST $p$ ($\pm$2\%)} \\
\midrule
All depths & 12,299 & $2.340 \pm 0.003$ & --- \\
Size-matched random & 3,951 & $2.370 \pm 0.003$ & ref \\
Depth 0+1 only & 3,951 & $2.390 \pm 0.003$ & 0.00012 \\
Binary-filtered & 3,951 & $2.393 \pm 0.003$ & 0.00009 \\
Depth 0 only & 797 & $2.685 \pm 0.004$ & $< 0.001$ \\
\bottomrule
\end{tabular}
\end{table}

At equal sample count ($n = 3{,}951$), random subsampling (ppl = 2.370) outperforms depth 0+1 filtering (ppl = 2.390) and binary-filtered selection (ppl = 2.393).
TOST confirms equivalence at the $\pm$2\% margin ($p = 0.00012$) but not at $\pm$1\% ($p = 0.087$), indicating practical equivalence with limited statistical precision due to only 3~random seeds.
Depth-0-only training yields substantially higher perplexity (2.685), but this reflects insufficient data volume (797~samples) rather than a depth quality effect.

This negative finding is robust across seeds (within-variant variance $< 0.003$).
Depth fingerprinting should therefore be positioned as a \emph{diagnostic} tool for provenance characterization---determining how much recursive generation a corpus contains---rather than a \emph{curative} mechanism for data selection.

%% ============================================================
%% 5. ANALYSIS
%% ============================================================
\section{Analysis}

\subsection{Reference Model Robustness}

To verify that DDR is not an artifact of the Pythia-410M reference model, we recompute features using OLMo-2-1B and Qwen3-0.6B as alternative references (Table~\ref{tab:refmodel}).

\begin{table}[t]
\centering
\small
\caption{DDR stability across reference models. Drift direction is consistent; $|\rho|$ magnitude varies moderately.}
\label{tab:refmodel}
\begin{tabular}{lccc}
\toprule
\textbf{Reference Model} & \textbf{LOGO F1} & \textbf{ent.\ mean $\rho$} & \textbf{tail decay $\rho$} \\
\midrule
Pythia-410M & 0.788 & $-$0.52 & $+$0.05 \\
OLMo-2-1B & 0.780 & $-$0.44 & $+$0.07 \\
Qwen3-0.6B & 0.795 & $-$0.48 & $+$0.04 \\
\bottomrule
\end{tabular}
\end{table}

All three references yield consistent drift direction, confirming DDR is not Pythia-specific.
LOGO F1 varies by $< 1.5$~pp across references, indicating classifier performance is robust to reference choice.
However, the specific $\rho$ values for entropy mean fluctuate by $\pm 0.08$ (from $-0.44$ to $-0.52$), reinforcing the classification of entropy mean as moderate rather than strong drift.
Tail decay and entropy max remain near-stationary across all references ($|\rho| < 0.10$), confirming that shape-feature stability is a genuine structural property of DDR.

\subsection{Prompt-Type Dependence}

We evaluate three prompt types---continuation, inspired rewriting, and elaboration---using Pythia-2.8B with 1,000 seeds per type (Table~\ref{tab:prompt}).

\begin{table}[t]
\centering
\small
\caption{Prompt-type comparison. Continuation chains yield the strongest depth signal; differences are moderate.}
\label{tab:prompt}
\begin{tabular}{lcccc}
\toprule
\textbf{Prompt Type} & \textbf{CV F1} & \textbf{$d_4/d_0$} & \textbf{hapax $\rho$} & \textbf{ent.\ mean $\rho$} \\
\midrule
Continuation & 0.773 & 0.835 & $-$0.72 & $-$0.50 \\
Rewriting & 0.733 & 0.938 & $-$0.78 & $-$0.45 \\
Elaboration & 0.723 & 0.984 & $-$0.81 & $-$0.43 \\
\bottomrule
\end{tabular}
\end{table}

Continuation chains achieve the highest classifier F1 (0.773) with the lowest depth retention ratio ($d_4/d_0 = 0.835$), indicating that direct continuation maximizes cumulative lexical compression.
Elaboration chains retain the highest depth retention (0.984) but yield the weakest classifier signal (0.723), reflecting a tradeoff between semantic divergence and depth discriminability.
The inter-prompt-type gap is $\sim$5~pp, indicating moderate prompt-type dependence.
Leave-one-prompt-type-out (LOPTO) evaluation yields an average F1 of 0.748, confirming reasonable cross-prompt transfer.

%% ============================================================
%% 6. DISCUSSION
%% ============================================================
\section{Discussion and Limitations}

\paragraph{Three-gradient DDR.}
The three-gradient structure (strong / moderate / near-stationary) provides a more nuanced characterization than a binary drifting/stationary split.
This structure is interpretable: strong-drift features capture vocabulary impoverishment (a cumulative physical process), moderate-drift features capture distributional location shifts (attenuated by model capacity and architecture), and near-stationary features reflect preserved distributional shape.
The gradient structure is robust across reference models and domains but scale-dependent---at 8B parameters, entropy mean approaches near-stationary territory.

\paragraph{Practical significance of the 7-feature advantage.}
The 2.2~pp margin over the best single-detector baseline (avg-logprob) is statistically significant ($p = 0.008$) but practically modest.
We position the 7-feature ensemble on interpretability and efficiency rather than absolute performance superiority.
The feature-level DDR analysis provides diagnostic value beyond the classifier itself: it reveals \emph{which} distributional properties are affected by recursive generation and at what rate.

\paragraph{Limitations.}
\textbf{(1) Base models only:} Our experiments use base (pre-trained) models. Instruction-tuned models impose length regulation and format constraints that compress the distributional variation DDR exploits.
\textbf{(2) Temperature sensitivity:} All experiments use $t = 0.7$; at $t = 0.9$, the signal degrades substantially, restricting the reliable operating range to $t \in [0.5, 0.7]$.
\textbf{(3) Scale ceiling:} Entropy-based features degrade at 8B; frontier-scale models ($> 14$B) may require surface-only classification, which limits resolution.
\textbf{(4) 3-class resolution limit:} 5-class ordinal classification yields F1 $\approx 0.495$ with no improvement from ordinal methods (Appendix~\ref{app:ordinal}); gen-1/gen-2+ is the practical resolution boundary.
\textbf{(5) Single perplexity metric:} The downstream experiment uses only CLM perplexity with 3~seeds, providing limited evidence for downstream effects.

%% ============================================================
%% 7. CONCLUSION
%% ============================================================
\section{Conclusion}

We introduced generation depth classification as a task beyond binary detection and characterized differential depth response (DDR) across 5 generators from 3 model families.
A lightweight 7-feature ensemble achieves LOGO binary F1 of 0.675 (range 0.620--0.715 per generator) with a single 410M-parameter reference model, outperforming single-detector baselines by 2.2~pp.
DDR follows a three-gradient structure---strong drift in surface features, moderate drift in entropy features, near-stationarity in shape features---that is robust across reference models and domains but degrades at larger model scales.
Depth-aware data curation is statistically equivalent to random subsampling, positioning depth fingerprinting as a diagnostic tool for provenance characterization.
Future work should validate DDR at $> 14$B scale and investigate whether architecture-specific feature engineering can recover the signal lost at larger model capacities.

%% ============================================================
%% REFERENCES
%% ============================================================
\begin{thebibliography}{20}

\bibitem[Bao et~al., 2024]{bao2024fastdetectgpt}
Guangsheng Bao, Yanbin Zhao, Zhiyang Teng, Linyi Yang, and Yue Zhang. 2024.
\newblock Fast-DetectGPT: Efficient zero-shot detection of machine-generated text via conditional probability curvature.
\newblock In \emph{Proc.\ ICLR}.

\bibitem[Briesch et~al., 2023]{briesch2023large}
Martin Briesch, Dominik Sobania, and Franz Rothlauf. 2023.
\newblock Large language models suffer from their own output: An analysis of the self-consuming generative process.
\newblock \emph{arXiv:2311.16822}.

\bibitem[Cao et~al., 2020]{cao2020rank}
Wenzhi Cao, Vahid Mirjalili, and Sebastian Raschka. 2020.
\newblock Rank consistent ordinal regression for neural networks with CORAL.
\newblock \emph{Pattern Recognition}, 108:107450.

\bibitem[Chen et~al., 2025]{chen2025llm}
Yun-Da Chen, Yen-Ting Lin, and Yun-Nung Chen. 2025.
\newblock LLM as a broken telephone: Iterative generation introduces attractor cycles.
\newblock In \emph{Proc.\ ACL}.

\bibitem[Chen and Guestrin, 2016]{chen2016xgboost}
Tianqi Chen and Carlos Guestrin. 2016.
\newblock XGBoost: A scalable tree boosting system.
\newblock In \emph{Proc.\ KDD}, pages 785--794.

\bibitem[Gehrmann et~al., 2019]{gehrmann2019gltr}
Sebastian Gehrmann, Hendrik Strobelt, and Alexander Rush. 2019.
\newblock GLTR: Statistical detection and visualization of generated text.
\newblock In \emph{Proc.\ ACL: System Demonstrations}.

\bibitem[Guo et~al., 2024]{guo2024curious}
Zi~Guo, Wen-Sen Lu, and Wei~Xu. 2024.
\newblock Curious repetitions: The compressing nature of recursive LLM generation.
\newblock In \emph{Proc.\ EMNLP}.

\bibitem[Hans et~al., 2024]{hans2024binoculars}
Abhimanyu Hans, Avi Schwarzschild, Valeriia Cherepanova, Hamid Kazemi, Aniruddha Saha, Micah Goldblum, Jonas Geiping, and Tom Goldstein. 2024.
\newblock Spotting LLMs with Binoculars: Zero-shot detection of machine-generated text.
\newblock In \emph{Proc.\ ICML}.

\bibitem[Kirchenbauer et~al., 2023]{kirchenbauer2023watermark}
John Kirchenbauer, Jonas Geiping, Yuxin Wen, Jonathan Katz, Ian Miers, and Tom Goldstein. 2023.
\newblock A watermark for large language models.
\newblock In \emph{Proc.\ ICML}.

\bibitem[Liang et~al., 2024]{liang2024mapping}
Weixin Liang, Yaohui Zhang, Zhengxuan Wu, Haley Lepp, Wenlong Ji, Xuandong Zhao, Hancheng Cao, Sheng Liu, Siyu He, Zhi Huang, et~al. 2024.
\newblock Mapping the increasing use of LLMs in scientific papers.
\newblock \emph{arXiv:2404.01268}.

\bibitem[Mitchell et~al., 2023]{mitchell2023detectgpt}
Eric Mitchell, Yoonho Lee, Alexander Khazatsky, Christopher Manning, and Chelsea Finn. 2023.
\newblock DetectGPT: Zero-shot machine-generated text detection using probability curvature.
\newblock In \emph{Proc.\ ICML}.

\bibitem[Penedo et~al., 2024]{penedo2024fineweb}
Guilherme Penedo, Hynek Kydl{\'\i}{\v{c}}ek, Loubna Ben~Allal, Anton Lozhkov, Margaret Mitchell, Colin Raffel, Leandro Von Werra, and Thomas Wolf. 2024.
\newblock The FineWeb datasets: Decanting the web for the finest text data at scale.
\newblock In \emph{Proc.\ NeurIPS Datasets and Benchmarks}.

\bibitem[Shumailov et~al., 2024]{shumailov2024model}
Ilia Shumailov, Zakhar Shumaylov, Yiren Zhao, Nicolas Papernot, Ross Anderson, and Yarin Gal. 2024.
\newblock AI models collapse when trained on recursively generated data.
\newblock \emph{Nature}, 631:755--759.

\bibitem[Yang et~al., 2024]{yang2024dna}
Xianjun Yang, Wei Cheng, Yue Wu, Linda Petzold, William Yang Wang, and Haifeng Chen. 2024.
\newblock DNA-GPT: Divergent n-gram analysis for training-free detection of GPT-generated text.
\newblock In \emph{Proc.\ ICLR}.

\end{thebibliography}

%% ============================================================
%% APPENDIX
%% ============================================================
\clearpage
\appendix

\section{Finer-Grained Ordinal Classification}
\label{app:ordinal}

We evaluate 5-class ordinal classification (gen-0 through gen-4) using standard XGBoost, CORAL \cite{cao2020rank}, and Frank--Hall ordinal decomposition (Table~\ref{tab:ordinal}).

\begin{table}[h]
\centering
\small
\caption{Ordinal granularity comparison. 5-class classification collapses; ordinal methods provide no improvement.}
\label{tab:ordinal}
\begin{tabular}{lccc}
\toprule
\textbf{Method} & \textbf{5-class F1} & \textbf{3-class F1} & \textbf{$\Delta$F1} \\
\midrule
XGBoost (standard) & 0.495 & 0.788 & $-$0.293 \\
CORAL & 0.498 & --- & --- \\
Frank--Hall & 0.492 & --- & --- \\
\bottomrule
\end{tabular}
\end{table}

The gen-1/gen-2 boundary is the primary bottleneck (43\% confusion rate), worsening at deeper boundaries: gen-2/gen-3 (52\%) and gen-3/gen-4 (58\%).
Binary classification errors concentrate as false negatives (85\% of errors): gen-2+ texts are misclassified as gen-1 rather than vice versa.
The three-class partition (gen-0 / gen-1 / gen-2+) represents the practical resolution ceiling for our feature set.

\section{Human Edit Robustness}
\label{app:human_edit}

We evaluate robustness to simulated human editing through random word replacement and LLM-simulated rephrasing (Table~\ref{tab:edit}).

\begin{table}[h]
\centering
\small
\caption{Robustness to simulated editing. Signal degrades gracefully under light to moderate perturbation.}
\label{tab:edit}
\begin{tabular}{lccc}
\toprule
\textbf{Edit Type} & \textbf{Binary F1} & \textbf{Drop (pp)} & \textbf{95\% CI} \\
\midrule
None (baseline) & 0.680 & --- & {[0.663, 0.697]} \\
Random 5\% & 0.636 & $-$4.4 & {[0.618, 0.654]} \\
Random 10\% & 0.602 & $-$7.8 & {[0.583, 0.621]} \\
Random 20\% & 0.575 & $-$10.5 & {[0.556, 0.594]} \\
LLM rephrase & 0.632 & $-$4.8 & {[0.612, 0.652]} \\
\bottomrule
\end{tabular}
\end{table}

LLM rephrasing (affecting $\sim$50\% of tokens by edit distance) causes a 4.8~pp drop, comparable to 5\% random word replacement ($-$4.4~pp).
The signal remains above chance (0.50) even at 20\% random replacement (F1 = 0.575), indicating that depth fingerprints, while degraded, are not entirely erased by moderate editing.
The relatively gentle impact of LLM rephrasing (compared to its high edit distance) reflects that synonym substitution preserves more distributional structure than random replacement.

\section{LD-Score as Additional Feature}
\label{app:ld_score}

We test LD-Score---a character-level distribution signature---as an eighth feature (Table~\ref{tab:ldscore}).

\begin{table}[h]
\centering
\small
\caption{LD-Score provides minimal marginal gain due to redundancy with hapax ratio.}
\label{tab:ldscore}
\begin{tabular}{lcc}
\toprule
\textbf{Feature Set} & \textbf{CV F1} & \textbf{$\Delta$ vs 7-feat} \\
\midrule
7 features & 0.788 & --- \\
7 features + LD-Score & 0.798 & $+$1.0~pp \\
LD-Score alone & 0.682 & $-$10.6~pp \\
\bottomrule
\end{tabular}
\end{table}

LD-Score alone achieves ordinal $\rho = {-}0.42$ with depth (moderate drift), but correlates strongly with hapax ratio ($r_s = 0.72$), making its marginal contribution redundant.
The +1.0~pp gain does not justify increased feature complexity.
However, the existence of character-level depth drift provides supplementary evidence that DDR operates across multiple linguistic granularities.

\section{End-to-End Classifier Comparison}
\label{app:e2e}

We compare the 7-feature XGBoost model against a QLoRA-fine-tuned Qwen3-0.6B end-to-end classifier (Table~\ref{tab:e2e}).

\begin{table}[h]
\centering
\small
\caption{End-to-end vs.\ handcrafted features. QLoRA excels within distribution but fails to generalize under LOGO.}
\label{tab:e2e}
\begin{tabular}{lcccc}
\toprule
\textbf{Method} & \textbf{CV 3-cl.} & \textbf{CV Bin.} & \textbf{LOGO Bin.} & \textbf{Cost} \\
\midrule
7-feature XGBoost & 0.788 & 0.680 & 0.675 & 1$\times$ \\
QLoRA Qwen3-0.6B & 0.882 & 0.840 & 0.620 & 12.4$\times$ \\
\bottomrule
\end{tabular}
\end{table}

QLoRA achieves substantially higher within-distribution performance (+9.4~pp for 3-class, +16~pp for binary) by learning fine-grained distributional patterns beyond our 7 handcrafted features.
However, its LOGO binary F1 (0.620) falls \emph{below} the 7-feature model (0.675), revealing severe overfitting to generator-specific artifacts.
The LOGO-vs-CV gap is 22~pp for QLoRA vs.\ 10~pp for 7-feature XGBoost, and inference requires 8.5~GB GPU memory at 12.4$\times$ the computational cost.
The 7-feature approach is positioned on its generalization and efficiency advantages.

\section{Text Length Confound Analysis}
\label{app:length}

We evaluate whether text length confounds the depth signal (Table~\ref{tab:length}).

\begin{table}[h]
\centering
\small
\caption{Length is not a significant confound. Depth signal persists after length normalization and matching.}
\label{tab:length}
\begin{tabular}{lcc}
\toprule
\textbf{Metric} & \textbf{Value} & \textbf{Conclusion} \\
\midrule
Spearman $\rho$(length, depth) & 0.11 & Weak ($|\rho| < 0.15$) \\
Length-matched binary F1 & 0.652 & $-$2.8~pp from unmatched \\
Random subsample F1 (same $n$) & 0.656 & Similar drop $\to$ sample size effect \\
Normalized TTR drift preserved? & Yes & Depth signal independent of length \\
\bottomrule
\end{tabular}
\end{table}

The length--depth correlation is weak ($\rho = 0.11$, 95\% CI [$-$0.06, 0.14]).
Length-matched subsampling reduces F1 by 2.8~pp (0.680 $\to$ 0.652), but a random subsample of the same size shows a comparable drop (0.656), confirming this reflects sample-size reduction rather than confound removal.
Length-normalized TTR and hapax trends are preserved (and in fact strengthened), confirming that depth signal is independent of text length variation.

\section{Multi-Domain Evaluation}
\label{app:domain}

We evaluate depth fingerprinting across 4 domains: Wikipedia, news, creative writing, and academic text (Table~\ref{tab:domain}).

\begin{table}[h]
\centering
\small
\caption{Multi-domain binary F1. Depth signal is domain-robust; academic text shows the weakest signal.}
\label{tab:domain}
\begin{tabular}{lccc}
\toprule
\textbf{Domain} & \textbf{Binary F1} & \textbf{95\% CI} & \textbf{Seeds} \\
\midrule
Wikipedia & 0.692 & {[0.672, 0.712]} & 500 \\
News & 0.675 & {[0.655, 0.695]} & 500 \\
Creative & 0.660 & {[0.640, 0.680]} & 500 \\
Academic & 0.648 & {[0.628, 0.668]} & 500 \\
\midrule
Cross-domain LOGO & 0.665 & --- & --- \\
\bottomrule
\end{tabular}
\end{table}

The maximum inter-domain gap is 4.4~pp (Wikipedia vs.\ academic).
Academic text scores lowest because its inherently specialized vocabulary and high term reuse partially mask depth-induced diversity compression.
Cross-domain LOGO (train on 3 domains, test on held-out) yields F1 = 0.665, close to the overall average, confirming domain robustness.

\section{Paraphrastic Chain Comparison}
\label{app:paraphrastic}

We contrast paraphrastic chains (semantic-preserving rewrites) with continuation chains to probe the role of semantic constraints in depth discriminability (Table~\ref{tab:para}).

\begin{table}[h]
\centering
\small
\caption{Paraphrastic vs.\ continuation chains. Semantic preservation accelerates attractor convergence and suppresses depth signal.}
\label{tab:para}
\begin{tabular}{lcccc}
\toprule
\textbf{Chain Type} & \textbf{Bin.\ F1} & \textbf{hapax $\rho$} & \textbf{ent.\ mean $\rho$} & \textbf{Convergence} \\
\midrule
Continuation & 0.682 & $-$0.61 & $-$0.50 & $>$ depth 4 \\
Paraphrastic & 0.585 & $-$0.28 & $-$0.15 & $\sim$ depth 2.8 \\
\bottomrule
\end{tabular}
\end{table}

Paraphrastic chains converge to attractor states by depth $\sim$2.8, consistent with the attractor cycle theory of Chen et al.\ \cite{chen2025llm}.
Their binary F1 (0.585) remains above chance but approaches the boundary of practical utility.
Continuation chains, lacking semantic constraints, sustain depth signal beyond depth 4, confirming that non-paraphrastic generation is the core scenario for depth fingerprinting.
The 9.7~pp gap ($\Delta$ F1) quantifies the cost of semantic preservation for depth discriminability.

\end{document}
